% Consolidated 20-magnetostatics from 98655b03c6d59adeb7b2c93d6e27b4e1cd8378ea; original TeX fragments preserved.

% === BEGIN 20-magnetostatics-219-235.tex blob 39bf2eff627a4bf9183989e1e587fba0b2ffbecf ===
% Irodov Part III §3.5 problems 219--235. User source PDF pp.133--135,
% printed pages 124--126; figures 3.59--3.64 checked as page images.
\electromagproblem{219}{124}
半径 $R=100\,\mathrm{mm}$ 的细圆环通电流
$I=1.00\,\mathrm A$。求磁感应强度：
(a) 圆心；(b) 轴线上距圆心
$x=100\,\mathrm{mm}$ 处。
\exsolution
毕奥—萨伐尔定律给圆环轴向场
\[
B(x)=\frac{\mu_0IR^2}{2(R^2+x^2)^{3/2}}.
\]
(a) 令 $x=0$：
$\boxed{B_0=\mu_0I/(2R)\simeq6.3\,\mathrm{\mu T}}$。
(b) 令 $x=R$：
\[
\boxed{B(R)=\frac{\mu_0I}{2^{5/2}R}
\simeq2.2\,\mathrm{\mu T}\approx2.3\,\mathrm{\mu T}}.
\]
\excheck 方向按电流右手定则确定，
轴上两侧磁场方向相同。

\electromagproblem{220}{124}
电流 $I$ 沿内接于半径 $R$ 圆周的
正 $n$ 边形细导线流动。
求中心磁场并讨论 $n\to\infty$。
\exsolution
每条边距中心为
$a=R\cos(\pi/n)$，
两端相对中心垂足夹角均为 $\pi/n$。
有限直导线在中心贡献
\[
B_1=\frac{\mu_0I}{4\pi a}
[2\sin(\pi/n)]
=\frac{\mu_0I}{2\pi R}\tan\frac\pi n.
\]
$n$ 条边方向同向叠加：
\[
\boxed{B=\frac{\mu_0 nI}{2\pi R}\tan\frac\pi n}.
\]
因 $\lim_{n\to\infty}n\tan(\pi/n)=\pi$，
\[
\boxed{B\longrightarrow\frac{\mu_0I}{2R}}.
\]
\excheck 大边数极限与圆环中心磁场一致。

\electromagproblem{221}{124}
矩形电流框对角线长
$d=16\,\mathrm{cm}$，
两对角线夹角 $\phi=30^\circ$，
电流 $I=5.0\,\mathrm A$。
求中心磁场。
\exsolution
设矩形边长
$a=d\cos(\phi/2)$、
$b=d\sin(\phi/2)$。
由有限直导线公式，
两条长边与两条短边对中心的贡献
分别为
$B_{\rm long}=2\mu_0I a/(\pi bd)$、
$B_{\rm short}=2\mu_0I b/(\pi ad)$。
相加：
\[
B=\frac{2\mu_0I}{\pi d}
\left(\frac ab+\frac ba\right)
=\frac{2\mu_0I}{\pi d}
\frac{a^2+b^2}{ab}.
\]
因 $a^2+b^2=d^2$、
$2ab=d^2\sin\phi$，
\[
\boxed{B=\frac{4\mu_0I}{\pi d\sin\phi}
\simeq1.0\times10^{-4}\,\mathrm T=0.10\,\mathrm{mT}}.
\]
\excheck 对角线角度 $\phi$ 不是边与对角线的角；
若误代 $\phi/2$，系数会错误。

\electromagproblem{222}{124}
图3.59的细导线由半径 $R=120\,\mathrm{mm}$
的圆弧和底部直弦闭合，
弦对应圆心角 $2\phi=90^\circ$，
电流 $I=5.0\,\mathrm A$。
求中心 $O$ 处磁场。
\exsolution
图示圆弧所张圆心角为
$2\pi-2\phi$，
其磁场为
$B_{\rm arc}=\mu_0I(2\pi-2\phi)/(4\pi R)$。
底部弦距中心的垂直距离
$a=R\cos\phi$，
半弦长度为 $R\sin\phi$，
由直导线磁场公式
\[
B_{\rm chord}=
\frac{\mu_0I}{2\pi a}\sin\phi
=\frac{\mu_0I}{2\pi R}\tan\phi.
\]
依图示电流方向两者同向，
\[
\boxed{B=\frac{\mu_0I}{2\pi R}
(\pi-\phi+\tan\phi)}.
\]
$\phi=\pi/4$ 代入得到
$\boxed{B\simeq28\,\mathrm{\mu T}}$。
\excheck 不能把缺口角 $2\phi$
当作\emph{实际导电圆弧}的张角。

\electromagproblem{223}{124}
求图3.60(a)、(b)中电流为 $I$
的闭合导线在点 $O$ 处的磁场；
(a) 已知半径 $a,b$ 与角 $\phi$；
(b) 已知曲线半径 $a$ 和方形边 $b$。
\exsolution
(a) 分别对半径为 $a,b$ 的同心圆弧
和两段径向连线作毕奥—萨伐尔积分。
径向段有
$\mathrm d\bm \ell\parallel\bm r$
而贡献为零。
由图3.60(a)的顺时针大圆弧
角度 $2\pi-\phi$ 与相反方向小弧
角度 $\phi$，以图示合场方向为正：
\[
\boxed{B_a=\frac{\mu_0I}{4\pi}
\left(\frac{2\pi-\phi}{a}+\frac{\phi}{b}\right)}.
\]
(b) 图3.60(b)由半径 $a$
的 $3\pi/2$ 圆弧以及两段与 $O$ 垂线
距离为 $b$ 的直导线组成。
逐段积分给
$B_{\rm arc}=\mu_0I(3\pi/2)/(4\pi a)$
的相应图示圆弧贡献；结合图中弧段覆盖与
方形两直段到点 $O$ 的斜距与夹角（各自投影包含 $1/\sqrt2$），整理按原书图示为
\[
\boxed{B_b=\frac{\mu_0I}{4\pi}
\left(\frac{3\pi}{4a}+\frac{\sqrt2}{b}\right)}.
\]
\excheck 本题的最终系数强烈依赖图示
导线走向，须与原图3.60而不是文字“圆环”
核对；两小段径向线段的贡献严格为零。

\electromagproblem{224}{125}
长薄壁圆柱导电管半径 $R$，沿轴方向有一条
宽度 $h\ll R$ 的纵缝。总纵向电流 $I$
沿未切开的管壁均匀分布。
求管内距缝口 $r$ 处磁场的近似值。
\exsolution
若整圈电流片均匀，理想无限圆管内部场为零。
引入狭缝后，电流分布可视为整圈电流片
\emph{减去}宽度为 $h$ 的窄条电流，
窄条等效线电流模
$I_{\rm miss}\simeq Ih/(2\pi R)$。
在距狭缝 $r$ 大于其宽度、
但仍可将缺口视为细导线的区域，
毕奥—萨伐尔定律
\[
B\simeq\frac{\mu_0I_{\rm miss}}{2\pi r}
=\boxed{\frac{\mu_0Ih}{4\pi^2Rr}}.
\]
\excheck 此为 $h\ll R$ 且离狭缝不太近
的窄缝近似；$r\lesssim h$ 时
细线替代会失去精度。

\electromagproblem{225}{125}
图3.61：长直导线的横截面为半径 $R$
的细半圆形导电带，
沿轴向均匀通总电流 $I$。
求半圆中心 $O$ 的磁感应强度。
\exsolution
沿半圆弧角度 $\theta\in[0,\pi]$
分割并联的无限细轴向电流通道，
各通道电流 $\mathrm dI=(I/\pi)\mathrm d\theta$。
距离中心均为 $R$，
每条无限轴向细线产生
$\mathrm dB=\mu_0\,\mathrm dI/(2\pi R)$。
各磁场方向随 $\theta$ 变化，
横向分量相消，沿半圆对称轴的分量相加：
\[
B=\frac{\mu_0I}{2\pi^2R}
\int_0^\pi\sin\theta\,\mathrm d\theta
=\boxed{\frac{\mu_0I}{\pi^2R}}.
\]
\excheck 非无限电流线\emph{沿半圆弧运行}，
而是纵向电流分布在半圆形横截面上；
两种几何给出的系数不同。

\electromagproblem{226}{125}
图3.62(a)、(b)、(c)中电流为 $I$，
弯曲段半径为 $R$，直线段可认为无限长。
分别求各构型圆心 $O$ 的磁感应强度。
\exsolution
利用圆弧磁场
$B_{\rm arc}=\mu_0I\Delta\theta/(4\pi R)$，
及有限直线磁场
$B_{\rm line}=\mu_0I(\sin\theta_1+
\sin\theta_2)/(4\pi d)$。
(a) 图中直线与观察点的连线共线，
直线部分 $\mathrm d\bm\ell\times\bm r=0$，
仅半圆弧贡献：
\[
\boxed{B_a=\frac{\mu_0I}{4\pi}\frac{\pi}{R}
=\frac{\mu_0I}{4R}}.
\]
(b) 图中圆弧及两半无限直线贡献相加：
\[
\boxed{B_b=\frac{\mu_0I}{4\pi R}
\left(1+\frac{3\pi}{2}\right)}.
\]
(c) 半圆/长圆弧与两半无限直线同向叠加：
\[
\boxed{B_c=\frac{\mu_0I}{4\pi R}
(2+\pi)}.
\]
\excheck 三种导线布局不同，
需按原图3.62计入各段；
这里列出的量均为 $O$ 处模，
正负方向由图示电流右手规则确定。

\electromagproblem{227}{125}
一根很长的细直导线在一点折成直角，
通电流 $I=5.0\,\mathrm A$。
在折角处沿垂直导线所在平面的方向距离
$\ell=35\,\mathrm{cm}$ 处求磁场。
\exsolution
每条半无限导线到观察点的垂距为 $\ell$，
由毕奥—萨伐尔积分得到
磁场模
$B_1=B_2=\mu_0I/(4\pi\ell)$。
两段导线方向互相垂直，
该两段贡献也互相垂直，
所以
\[
\boxed{B=\sqrt{B_1^2+B_2^2}
=\frac{\sqrt2\,\mu_0I}{4\pi\ell}}.
\]
代入：
$\boxed{B\simeq2.0\,\mathrm{\mu T}}$。
\excheck 不是两条同向磁场简单相加成2倍，
应作矢量合成得 $\sqrt2$。

\electromagproblem{228}{125}
图3.63(a)、(b)、(c)，
半径 $R=100\,\mathrm{mm}$ 的圆形弯段
与数条长直导线连接，通电流
$I=8.0\,\mathrm A$。
求图中点 $O$ 的磁感应强度。
\exsolution
将图中各弯段、空间直导线按毕奥—萨伐尔定律
逐段求 $\bm B$，再在
$\bm e_x,\bm e_y,\bm e_z$ 上矢量合成。
对每段轴向或径向直线，
若 $\mathrm d\bm\ell\parallel\bm r$，
其场为零；对半无限直线采用
$\mu_0I/(4\pi R)$，
对整圆弧采用 $\mu_0I/(2R)$。
根据图3.63各方向整理，三种模分别为
\[
\boxed{B_a=\frac{\mu_0I}{4\pi R}\sqrt{4+\pi^2}},
\quad
\boxed{B_b=\frac{\mu_0I}{4\pi R}
\sqrt{2+2\pi+\pi^2}},
\]
\[
\boxed{B_c=\frac{\mu_0I}{4\pi R}\sqrt2}.
\]
代入给定电流与长度：
\[
\boxed{B_a\simeq30\,\mathrm{\mu T},\quad
B_b\simeq34\,\mathrm{\mu T},\quad
B_c\simeq11\,\mathrm{\mu T}}.
\]
\excheck \textbf{原书书末的数值单位疑误：}
本版 PDF 所印磁场模似写
0.30、0.34、0.11 $\mathrm{\mu T}$，
但直接把书末\emph{自身公式}
代入 $I=8$A、$R=0.10$m
得30、34、11 $\mathrm{\mu T}$，
相差100倍；保留原公式并记录数值冲突。

\electromagproblem{229}{125}
(a) 无限大平面沿面内有均匀表面电流密度
$\bm K$（单位 $\mathrm{A/m}$），
求两侧磁场模与方向；
(b) 两平行平面分别有 $\bm K,-\bm K$，
求各区域磁场。
\exsolution
对单个无限电流片，
取横跨片的长方形安培环路，
两侧切向磁场大小相等反向，
\[
2B\ell=\mu_0K\ell,
\quad\boxed{B=\frac{\mu_0K}{2}}.
\]
方向由
$\bm B_\pm=\pm\frac{\mu_0}{2}
(\bm K\times\bm n)$
依所取法线约定确定。
(b) 两片电流反向，片间磁场同向叠加，
片外磁场相消，
\[
\boxed{B_{\rm between}=\mu_0K,\qquad
B_{\rm outside}=0}.
\]
\excheck $\bm K$ 是\emph{单位长度电流}，
不能误认为每单位面积的体电流密度 $j$。

\electromagproblem{230}{125--126}
无限厚平板占据
$-d\le x\le d$，
内部流动平行于板面的恒定体电流密度
$j$，周围磁导率为真空值。
求磁场模随到中面距离 $|x|$ 的变化。
\exsolution
沿电流和法向构造安培长方环路。
由对称性板中面 $B(0)=0$，
板内穿过的电流每单位边长为
$2jx$。安培环路两长边贡献
$2B(x)\ell=\mu_0(2jx\ell)$，
故
\[
\boxed{B(x)=
\begin{cases}
\mu_0j|x|,& |x|\le d,\\
\mu_0jd,& |x|\ge d.
\end{cases}}
\]
磁场方向在两侧相反，
可用带符号表达
$B_y(x)=\mu_0jx$。
\excheck 在 $x=\pm d$ 连续；
板外不是随距离继续增大，
而是恒定磁场。

\electromagproblem{231}{126}
图3.64：长直导线上的直流电
$I$ 沿竖直方向流到点 $O$，
然后沿垂直于导线的无限导电平面
径向向外均匀扩散。
求空间中各处磁感应强度。
\exsolution
取圆柱坐标，以竖直导线为轴、
平面为 $z=0$。
电流和几何关于轴旋转对称，
$\bm B$ 仅可能沿 $\bm e_\varphi$。
在导线所在的半空间，
半径 $r$ 的圆形安培环路包围纵向电流
$I$，因此
\[
\boxed{B_\varphi(r,z>0)=
\frac{\mu_0I}{2\pi r}}.
\]
平面另一侧的安培环路不包围
穿过所张圆盘的电流，且无限径向电流
对称叠加给
\[
\boxed{B_\varphi(r,z<0)=0}.
\]
\excheck 边界 $z=0$ 的表面径向电流
允许磁场切向分量跃变；
不必让两侧磁场连续。

\electromagproblem{232}{126}
圆形电流环通电流 $I$，
计算其轴线上从 $-\infty$ 到
$+\infty$ 的线积分
$\int\bm B\cdot\mathrm d\bm r$，
并解释结果。
\exsolution
圆环半径记 $R$，
轴上磁感应强度
$B_z(z)=\mu_0IR^2/[2(R^2+z^2)^{3/2}]$。
因此
\[
\int_{-\infty}^{\infty}B_z\,\mathrm dz
=\frac{\mu_0IR^2}{2}
\left[\frac{z}{R^2\sqrt{R^2+z^2}}\right]_{-\infty}^{\infty}
=\boxed{\mu_0I}.
\]
另一说明：在无穷远处把轴线的两端
通过一段不穿过电流的远方路径闭合，
远方磁场的线积分趋零，
则安培环路定律给 $\mu_0I$。
\excheck 积分结果不依赖圆环半径；
此为有电流环穿链的拓扑性质。

\electromagproblem{233}{126}
一根无限长均匀圆截面导线半径 $R$，
通均匀直流体电流密度矢量 $\bm j$。
求距轴的横向位置矢量为 $\bm r$
处的磁感应强度矢量，内外均为真空磁导率。
\exsolution
轴对称圆柱导线
总电流 $I=j\pi R^2$，
圆周安培环路给
$B(2\pi r)=\mu_0I_{\rm enc}$。
在 $r\le R$，包围电流
$I_{\rm enc}=j\pi r^2$；
在 $r\ge R$，为
$j\pi R^2$。
用右手定则写矢量：
\[
\boxed{\bm B(\bm r)=
\begin{cases}
\dfrac{\mu_0}{2}\,\bm j\times\bm r,
&r\le R,\\[2mm]
\dfrac{\mu_0R^2}{2r^2}\,
\bm j\times\bm r,&r\ge R.
\end{cases}}
\]
\excheck 表面 $r=R$ 内外连续；
$\bm j$ 与 $\bm r$ 垂直，
磁场沿方位角方向。

\electromagproblem{234}{126}
一根长均匀圆截面直导线内，
挖出一条与原轴平行的长圆柱空腔，
两轴位移矢量为 $\bm a$，
大小为 $\ell$。
其余材料通均匀电流密度 $\bm j$。
求空腔内部磁场，特别讨论 $\ell=0$。
\exsolution
空腔可视为一根\emph{完整}均匀导电圆柱
和一根位于空腔位置、密度为 $-\bm j$
的圆柱叠加。
圆柱内部
$\bm B=\frac{\mu_0}{2}\bm j\times\bm r$。
在空腔内部某点，
相对于原轴是 $\bm r$，
相对于腔轴是 $\bm r-\bm a$，
所以
\[
\bm B_{\rm cavity}=
\frac{\mu_0}{2}
[\bm j\times\bm r-\bm j\times(\bm r-\bm a)]
=\boxed{\frac{\mu_0}{2}\bm j\times\bm a}.
\]
故空腔中磁场均匀，
大小为 $\mu_0j\ell/2$，
若 $\ell=0$：
$\boxed{\bm B=0}$。
\excheck 与均匀带电圆柱偏心空腔
的叠加技巧相似；
方向按 $\bm j\times\bm a$。

\electromagproblem{235}{126}
径向对称、沿轴向运动的电子束内，
磁感应强度模满足
$B(r)=br^\alpha$，
$b,\alpha>0$。
求电子束的轴向电流密度模 $j(r)$。
\exsolution
对以束轴为中心的圆形安培环路：
\[
2\pi r B(r)=\mu_0 I_{\rm enc}(r).
\]
而
$I_{\rm enc}(r)=2\pi\int_0^rj(s)s\,\mathrm ds$。
因此
\[
j(r)=\frac1{\mu_0r}
\frac{\mathrm d}{\mathrm dr}[rB(r)]
=\frac1{\mu_0r}\frac{\mathrm d}{\mathrm dr}
(br^{\alpha+1}).
\]
得到
\[
\boxed{j(r)=\frac{b(\alpha+1)}{\mu_0}
r^{\alpha-1}}.
\]
\excheck 本题求大小；
真实电子运动方向与传统电流方向相反。
% === END 20-magnetostatics-219-235.tex ===

% === BEGIN 20-magnetostatics-236-250.tex blob 9ba5e102ce2306b48ec277252526df54d38af90a ===
% Irodov Part III §3.5 236--250, source PDF pp.135--137
% printed 126--128, original figures 3.65--3.67 reviewed.
\electromagproblem{236}{126}
有限长单层螺线管，轴向长度 $\ell$，
截面半径 $R$，单位长度匝数 $n$，
通电流 $I$。求其中心的磁感应强度。
\exsolution
轴上单个圆环距观察点 $z$ 时产生
$\mathrm dB=\mu_0 IR^2/
[2(R^2+z^2)^{3/2}]\,n\,\mathrm dz$。
中心处将各匝从
$-\ell/2$ 积分至 $+\ell/2$：
\[
B=\frac{\mu_0nIR^2}{2}
\int_{-\ell/2}^{\ell/2}
\frac{\mathrm dz}{(R^2+z^2)^{3/2}}.
\]
用积分原函数
$z/[R^2\sqrt{R^2+z^2}]$，
\[
\boxed{B=
\frac{\mu_0nI}{\sqrt{1+(2R/\ell)^2}}}.
\]
\excheck $\ell\gg R$ 时趋于
$\mu_0nI$；有限长度的场小于无限螺线管极限。

\electromagproblem{237}{126}
很长直螺线管，半径 $R$，
单位长度匝数 $n$，电流 $I$。
以端部为 $x=0$，
$x>0$ 沿轴指向管外，
$x<0$ 指向管内。
求 (a) 轴向场 $B(x)$ 与形状；
(b) 管内场离中部值
$B_0=\mu_0nI$ 相差 $\eta=1\%$
的位置至端部距离。
\exsolution
把螺线管的环匝看成
占据 $z\le0$ 的连续分布，
由毕奥—萨伐尔定律
\[
B(x)=\frac{\mu_0nIR^2}{2}
\int_{-\infty}^{0}
\frac{\mathrm dz}{[R^2+(x-z)^2]^{3/2}}.
\]
令 $u=x-z$ 后积分，得
\[
\boxed{B(x)=\frac{\mu_0nI}{2}
\left(1-\frac{x}{\sqrt{x^2+R^2}}\right)}.
\]
当 $x\to-\infty$，$B/B_0\to1$，
$x=0$ 时 $B/B_0=1/2$，
$x\to+\infty$ 时 $B\to0$。
(b) 令管内 $x=-x_0$，且
$B(-x_0)=B_0(1-\eta)$。
于是
$x_0/\sqrt{x_0^2+R^2}=1-2\eta$，
从而
\[
\boxed{x_0=
\frac{R(1-2\eta)}
{2\sqrt{\eta(1-\eta)}}
\simeq4.9R\approx5R}.
\]
\excheck 书末第237题直接给约 $5R$；
计算时必须选 $x<0$ 的\emph{管内}位置。

\electromagproblem{238}{126}
宽 $h=2.0\,\mathrm{cm}$ 的薄导电带
紧密螺旋绕成半径
$R=2.5\,\mathrm{cm}$ 的很长单层螺线管，
带内通电流 $I=5.0\,\mathrm A$。
求管内外磁场随离轴半径 $r$ 的变化。
\exsolution
薄带紧密盘绕，近似每轴向长度 $h$
对应一匝，绕向表面电流分量
$K_\varphi\simeq I/h$。
由无限长电流片磁场跳变条件，
绕向电流使管内产生
\[
\boxed{B_{\rm inside}\simeq
\frac{\mu_0 I}{h}\simeq0.31\,\mathrm{mT}},
\quad r<R.
\]
螺旋线还带一个很小的沿轴
流动的电流分量。
由每匝轴向进给 $h$、
周长 $2\pi R$，
有效轴向电流为
$I_z\simeq I h/(2\pi R)$，
故管外有方位向磁场
\[
\boxed{B_{\rm outside}(r)\simeq
\frac{\mu_0 I h}{4\pi^2Rr}},\quad r>R.
\]
更精细的几何需乘
$\cos\theta=
[1+(h/2\pi R)^2]^{-1/2}$
等微小螺旋倾角修正。
\excheck 本题的带是\emph{连续倾斜螺旋}，
不是一圈圈独立闭合圆环；
因此管外存在小的方位向磁场。

\electromagproblem{239}{126--127}
木质环形铁芯截面非常细，
均匀缠有 $N=2.5\times10^3$ 匝电流线圈，
每匝电流为 $I$。
求环形铁芯内部磁场与环形线圈
\emph{几何中心}磁场的比值。
\exsolution
设环形主半径为 $R$、截面半径很小。
环形铁芯内沿主圆周的环路，
安培定律得
$B_{\rm core}2\pi R=\mu_0NI$，
\[
B_{\rm core}\simeq\frac{\mu_0NI}{2\pi R}.
\]
环形线圈的外侧电流可用
等效中心环电流的首阶近似，
在环心产生
$B_{\rm centre}\simeq\mu_0I/(2R)$。
因此
\[
\boxed{\frac{B_{\rm core}}{B_{\rm centre}}
\simeq\frac N\pi
\simeq8.0\times10^2}.
\]
\excheck 薄截面假设允许忽略芯内径向
磁场非均匀性；木质材料磁导率近似为1。

\electromagproblem{240}{127}
长直均匀圆截面导线通电流
$I=10\,\mathrm A$。
求每单位轴向长度、贯穿导线
\emph{半个截面（从轴至表面的径向面）}
的磁通量。
\exsolution
设导线半径 $R$。
截面内的磁场
$B(r)=\mu_0Ir/(2\pi R^2)$，
方向沿方位角。
取包含导线轴的一条径向半平面，
其单位轴向长度的面积元为
$\mathrm dA=1\,\mathrm m\,\mathrm dr$，
面法线沿方位角，故
\[
\frac{\Phi}{L}=
\int_0^RB(r)\,\mathrm dr
=\frac{\mu_0I}{2\pi R^2}\frac{R^2}{2}
=\boxed{\frac{\mu_0I}{4\pi}}.
\]
代入
$\boxed{\Phi/L=1.0\,\mathrm{\mu Wb/m}}$。
\excheck 这里的“半个截面”是
\emph{沿轴伸展的径向半平面}；
若误取横截面所在的平面，
磁场切向该截面，通量必为零。

\electromagproblem{241}{127}
很长螺线管电流 $I$，
截面积 $S$，单位长度匝数 $n$。
求穿过\emph{端面平面}的磁通量。
\exsolution
远离端部的内部场
$B_0=\mu_0nI$。
端部轴向场近似为内部的一半，
$B_{\rm end}\simeq\mu_0nI/2$。
对端面面积 $S$ 积分，
在长螺线管、横截面尺度
相对螺线管长度很小时，
\[
\boxed{\Phi_{\rm end}\simeq
\frac12\mu_0nIS}.
\]
\excheck 与管内深处的截面磁通量
$\Phi_0=\mu_0nIS$ 比较，
端面约为其一半。

\electromagproblem{242}{127}
图3.65：矩形截面环形螺线管
总匝数 $N=1000$，电流
$I=1.7\,\mathrm A$，截面高
$h=5.0\,\mathrm{cm}$，
外径/内径 $\eta=1.6$。
求穿过环形芯截面的磁通量。
\exsolution
环形安培环路半径 $r$
上的磁场
$B(r)=\mu_0NI/(2\pi r)$。
截面从内半径 $a$ 到外半径 $b$，
$b/a=\eta$，
磁通量
\[
\Phi=\int_a^bB(r)h\,\mathrm dr
=\frac{\mu_0NIh}{2\pi}\ln\frac ba.
\]
因此
\[
\boxed{\Phi=
\frac{\mu_0NIh}{2\pi}\ln\eta}
=\boxed{8.0\,\mathrm{\mu Wb}}\quad\text{（约）}.
\]
\excheck 计算需要的是半径比，
外径比与内径比相同；
无需知道环形芯的绝对主半径。

\electromagproblem{243}{127}
半径为 $R=100\,\mathrm{mm}$
的细圆环，圆心磁感应强度
$B=6.0\,\mathrm{\mu T}$。
求圆环电流对应磁矩。
\exsolution
圆环圆心场
$B=\mu_0I/(2R)$，
由此求电流
$I=2RB/\mu_0$。
圆环磁矩
$p_m=I(\pi R^2)$，所以
\[
\boxed{p_m=\frac{2\pi R^3B}{\mu_0}}.
\]
代值：
\[
\boxed{p_m\simeq0.030\,\mathrm{A\,m^2}
=30\,\mathrm{mA\,m^2}}.
\]
\excheck 原书书末的纯文本提取可能把
$R^3$ 误写为 $R^2$；
由 $p_m=I\pi R^2$ 的量纲可以独立复核立方幂。

\electromagproblem{244}{127}
图3.66：$N=500$ 匝细绝缘导线
紧密绕在\emph{半个环面}上，
环面细圆截面直径
$d=5.0\,\mathrm{cm}$，
电流 $I=0.80\,\mathrm A$。
求整个绕组的磁矩。
\exsolution
每一小匝的面积约
$A=\pi(d/2)^2$，磁矩模为
$IA$，方向沿环面的切向。
半环面上的切线方向随环周
转过 $\pi$，对同一方向的投影积分：
\[
|\bm p_m|=
\frac{NIA}{\pi}
\left|\int_{-\pi/2}^{\pi/2}
\cos\theta\,\mathrm d\theta\right|
=\frac{2NIA}{\pi}.
\]
代入
$A=\pi d^2/4$，
得
\[
\boxed{p_m=\frac12NId^2
=\frac12(500)(0.8)(0.05)^2
\simeq0.50\,\mathrm{A\,m^2}}.
\]
\excheck 不能把所有小匝的磁矩当作平行，
否则会错用 $NIA$；
必须积分半环面上变化的磁矩方向。

\electromagproblem{245}{127}
图3.67：单层平面螺旋由
$N=100$ 匝紧密细线构成，
电流 $I=8\,\mathrm{mA}$，
最内、外匝半径分别
$a=50\,\mathrm{mm}$、$b=100\,\mathrm{mm}$。
求 (a) 螺旋中心磁感应强度；
(b) 总磁矩。
\exsolution
把匝数视为半径均匀连续分布，
$\mathrm dN=N\,\mathrm dr/(b-a)$。
每匝在中心贡献
$\mathrm dB=\mu_0I\,\mathrm dN/(2r)$，
故
\[
\boxed{B=
\frac{\mu_0IN}{2(b-a)}\ln\frac ba
\simeq7.0\,\mathrm{\mu T}}.
\]
磁矩为各环的
$I\pi r^2$ 叠加，
\[
p_m=I\pi\int_a^b r^2\,
\frac{N\,\mathrm dr}{b-a}
=\frac{\pi IN}{3(b-a)}(b^3-a^3).
\]
因 $b^3-a^3=(b-a)(a^2+ab+b^2)$：
\[
\boxed{p_m=
\frac{\pi IN}{3}(a^2+ab+b^2)
\simeq15\,\mathrm{mA\,m^2}}.
\]
\excheck 这是密匝连续近似；
与等间距离散100匝逐一相加只有小修正。

\electromagproblem{246}{127--128}
半径 $R$ 的薄绝缘圆盘，
一面均匀面带电密度 $\sigma$，
以角速度 $\omega$ 绕轴旋转。
求 (a) 圆心磁场；(b) 旋转电荷的磁矩。
\exsolution
半径 $r$ 宽 $\mathrm dr$
的环带带电
$\mathrm dq=2\pi\sigma r\,\mathrm dr$。
每旋转一周通过固定截面一次，
频率 $\omega/(2\pi)$，
等效环电流
$\mathrm dI=\sigma\omega r\,\mathrm dr$。
圆心处小圆环贡献
$\mathrm dB=\mu_0\mathrm dI/(2r)
=\mu_0\sigma\omega\,\mathrm dr/2$，
积分得
\[
\boxed{B(0)=\frac12\mu_0\sigma\omega R}.
\]
各小环磁矩
$\mathrm dp_m=(\pi r^2)\,\mathrm dI$，
故
\[
\boxed{p_m=\pi\sigma\omega\int_0^Rr^3\,\mathrm dr
=\frac14\pi\sigma\omega R^4}.
\]
\excheck $\sigma$ 与 $\omega$
的符号决定方向；
上述为正 $\sigma$ 的模与右手方向。

\electromagproblem{247}{128}
半径 $R=50\,\mathrm{mm}$ 的
不导电薄球壳，均匀面电荷密度
$\sigma=10.0\,\mathrm{\mu C/m^2}$，
绕过球心的轴以
$\omega=70\,\mathrm{rad/s}$ 旋转。
求球心磁感应强度。
\exsolution
用球面极角 $\theta$ 参数化，
半径 $r_\perp=R\sin\theta$
处的球面带电小环总电荷
$\mathrm dq=\sigma2\pi R^2\sin\theta\,\mathrm d\theta$。
等效小环电流
$\mathrm dI=\sigma\omega R^2\sin\theta\,\mathrm d\theta$。
该环到球心的距离都是 $R$，
球心轴向磁场分量为
\[
\mathrm dB_z=
\frac{\mu_0}{2}\frac{\mathrm dI\,(R\sin\theta)^2}
{R^3}
=\frac{\mu_0\sigma\omega R}{2}
\sin^3\theta\,\mathrm d\theta.
\]
利用
$\int_0^\pi\sin^3\theta\,\mathrm d\theta=4/3$，
\[
\boxed{B(0)=\frac23\mu_0\sigma\omega R
\simeq2.9\times10^{-11}\,\mathrm T=29\,\mathrm{pT}}.
\]
\excheck 不可把球壳当作平圆盘；
不同纬度圆环的轴向投影角不同。

\electromagproblem{248}{128}
质量 $m$、半径 $R$ 的均匀实心球
带总电荷 $q$，电荷在球内均匀分布。
绕穿过球心的轴以角速度 $\omega$ 旋转。
求磁矩 $p_m$ 及其与机械角动量 $L$ 的比。
\exsolution
均匀体电荷密度
$\rho_e=3q/(4\pi R^3)$。
绕轴距 $r_\perp$ 的体元
$\mathrm dq=\rho_e\mathrm dV$
转动产生磁矩微元
$\mathrm dp_m=\frac12\omega r_\perp^2\mathrm dq$，
因此
\[
p_m=\frac12\omega\rho_e
\int_{\rm sphere}r_\perp^2\,\mathrm dV
=\frac12\omega\rho_e\frac{8\pi R^5}{15}
=\boxed{\frac15 qR^2\omega}.
\]
均匀实心球转动惯量
$J=(2/5)mR^2$，
机械角动量 $L=J\omega$，
\[
\boxed{\frac{p_m}{L}=\frac{q}{2m}}.
\]
\excheck 磁矩可为带符号矢量，
若球带负电其磁矩方向与机械角动量相反。

\electromagproblem{249}{128}
无限长介质圆柱半径 $R$，
预先产生径向极化
$\bm P(r)=\alpha r\,\bm e_r$。
介质随后以角速度 $\omega$
绕圆柱轴旋转。
求圆柱中心磁场。
\exsolution
由极化的束缚电荷公式，
圆柱内束缚体密度
$\rho_b=-\nabla\cdot\bm P=-2\alpha$，
外侧面束缚面密度
$\sigma_b=P(R)=\alpha R$。
随材料转动，体电荷产生
$j_\varphi(r)=\rho_b\omega r=-2\alpha\omega r$，
表面电荷形成周向表面电流
$K_\varphi=\sigma_b\omega R=\alpha\omega R^2$。
无限长圆柱周向电流片外侧轴向磁场为零。
用安培边界条件，靠内表面场
$B_z(R^-)=\mu_0K_\varphi
=\mu_0\alpha\omega R^2$。
内部微分形式
$-\mathrm dB_z/\mathrm dr
=\mu_0j_\varphi=-2\mu_0\alpha\omega r$，
积分得
$B_z(r)=\mu_0\alpha\omega r^2$。
故球心（圆柱轴线）
\[
\boxed{B(0)=0}.
\]
\excheck 不能仅算带正束缚电荷的侧壁电流，
必须同时计入转动的负束缚\emph{体电荷}电流。

\electromagproblem{250}{128}
两质子以相同的平行速度
$v=300\,\mathrm{km/s}$ 运动。
求它们间磁相互作用力与静电力的比。
\exsolution
在非相对论近似下，
运动点电荷在另一质子处的磁场
$B\simeq(\mu_0/4\pi)qv/r^2$
（两速度与两质子连线垂直的情形）。
洛伦兹磁力
$F_m=qvB=\mu_0q^2v^2/(4\pi r^2)$，
静电力
$F_e=q^2/(4\pi\varepsilon_0r^2)$。
因此
\[
\boxed{\frac{F_m}{F_e}=
\mu_0\varepsilon_0v^2
=\left(\frac vc\right)^2}.
\]
代入 $v=3.0\times10^5\,\mathrm{m/s}$：
\[
\boxed{F_m/F_e\simeq1.0\times10^{-6}}.
\]
\excheck 对严格平行运动但质子相对位置
并不垂直速度时，磁力大小还包含
几何因子 $\sin^2\theta$。
原书答案采用两质子并列、连线横向的常见构型。
% === END 20-magnetostatics-236-250.tex ===

% === BEGIN 20-magnetostatics-251-262.tex blob 54e2f9a4f42817a40f80679fc1b7799cac800548 ===
% Irodov Part III §3.5 251--262, user PDF pp.137--139,
% printed pp.128--130; figures 3.68--3.73 visually reviewed.
\electromagproblem{251}{128}
图3.68中细导线通电流 $I=8.0\,\mathrm A$。
求图示点 $O$ 处单位长度导线所受磁力模与方向：
(a) 半圆弯曲半径 $R=10\,\mathrm{cm}$；
(b) U 形两长平行段距离 $\ell=20\,\mathrm{cm}$。
\exsolution
局部导线受力每单位长度
$\bm f=I\,\bm t\times\bm B_{\rm others}$。
(a) 半圆弧在其中心产生
$B_{\rm arc}=\mu_0I/(4R)$，
直径段中心处的其他平直部分磁场不贡献，
故
\[
\boxed{f_a=\frac{\mu_0I^2}{4R}
\simeq0.20\,\mathrm{mN/m}}.
\]
(b) 图中点 $O$ 在左竖直连接段上，
每条半无限水平长直线距离它 $\ell/2$，
两条电流的磁场方向在该位置相加：
$B=2\mu_0I/[4\pi(\ell/2)]
=\mu_0I/(\pi\ell)$，
所以
\[
\boxed{f_b=\frac{\mu_0I^2}{\pi\ell}
\simeq0.13\,\mathrm{mN/m}}.
\]
\excheck 对每段导线应用
$I\mathrm d\bm\ell\times\bm B$，
按图中的电流箭头确定力朝向弯曲内部或外部；
图示 $O$ 不等于整个回路的质心。

\electromagproblem{252}{128}
单层紧密铜线线圈半径 $R=30\,\mathrm{mm}$，
导线直径 $d=0.10\,\mathrm{mm}$，
电流 $I=10\,\mathrm{mA}$，
线圈轴与均匀外磁场方向重合。
估算使绕组拉断的外部磁感应强度。
\exsolution
圆环上的磁力每单位导线长度为
$f=IB$，近似沿径向向外。
圆环微小弧元两端张力的径向合力
$T\,\mathrm d\theta$
与外力
$IBR\,\mathrm d\theta$ 平衡，
所以张力 $T=IBR$。
若铜线的极限抗拉应力为 $\sigma_m$，
允许最大张力
$T_{\max}=\sigma_m A=
\sigma_m\pi d^2/4$。
断裂条件：
\[
\boxed{B_{\rm rupture}=
\frac{\pi d^2\sigma_m}{4RI}}.
\]
以铜线强度数量级
$\sigma_m\sim3\times10^8\,\mathrm{Pa}$
估算：
$\boxed{B_{\rm rupture}\sim8\times10^3\,\mathrm T
=8\,\mathrm{kT}}$。
\excheck 题面未给具体铜材的抗拉强度，
因此数值必须注明采用的材料强度假设；
理想估算的8 kT远超普通实验室磁场。

\electromagproblem{253}{128}
截面积 $S=2.5\,\mathrm{mm^2}$
的铜线弯成正方形三条边，
可绕水平轴 OO' 转动
（图3.69）。外磁场竖直，
通电流 $I=16\,\mathrm A$ 后偏转
$\theta=20^\circ$。
求磁场模。
\exsolution
设正方形边长为 $a$，
铜密度为 $\rho_{\rm Cu}$。
图示三条边的重力矩由重心几何求得：
\[
\tau_g=2\rho_{\rm Cu}Sa^2g\sin\theta.
\]
电流在底部与两侧形成的有效磁力矩为
$\tau_B=Ia^2B\cos\theta$。
平衡：
$Ia^2B\cos\theta
=2\rho_{\rm Cu}Sa^2g\sin\theta$。
因此
\[
\boxed{B=\frac{2\rho_{\rm Cu}gS}{I}\tan\theta}.
\]
取 $\rho_{\rm Cu}\simeq8.9\times10^3
\,\mathrm{kg/m^3}$，
\[
\boxed{B\simeq1.0\times10^{-2}\,\mathrm T
=10\,\mathrm{mT}}.
\]
\excheck 两边含 $a^2$ 的矩相消，
故题目不需给正方形边长。

\electromagproblem{254}{128--129}
图3.70：$N=200$ 匝小线圈
位于磁极间，面积
$S=1.0\,\mathrm{cm^2}$，
以一端悬于天平横杆上，
力臂 $\ell=30\,\mathrm{cm}$。
线圈电流
$I=22\,\mathrm{mA}$。
追加配重
$\Delta m=60\,\mathrm{mg}$
才恢复平衡。
求该处磁感应强度。
\exsolution
线圈磁矩
$p_m=NIS$。
外磁场对它产生力矩模
$\tau_B=NISB$（图中转角使磁矩与场方向垂直）。
追加配重产生的重力矩
$\tau_g=\Delta mg\ell$。
平衡条件
$NISB=\Delta mg\ell$，
于是
\[
\boxed{B=\frac{\Delta mg\ell}{NIS}}.
\]
SI 单位：
$\Delta m=6.0\times10^{-5}\,\mathrm{kg}$，
$S=10^{-4}\,\mathrm{m^2}$，
\[
\boxed{B\simeq0.40\,\mathrm T}.
\]
\excheck 重量单位是毫克 mg，
不是克 g；换算错误会放大1000倍。

\electromagproblem{255}{129}
边长 $a=8.0\,\mathrm{cm}$ 的正方形电流框
$I=0.90\,\mathrm A$，
与无限长直导线电流
$I_0=5.0\,\mathrm A$ 同处平面。
电流框的转轴与长导线平行，
轴距是 $\eta a$，$\eta=1.5$。
求 (a) 合安培力；(b) 恒定电流下
绕该轴翻转 $180^\circ$ 的外功。
\exsolution
长线在距线 $r$ 处
$B(r)=\mu_0I_0/(2\pi r)$。
电流框的两条长边距直线分别为
$r_1=a(\eta-\tfrac12)$、
$r_2=a(\eta+\tfrac12)$。
两边受力方向相反，
其差为
\[
F=Ia[B(r_1)-B(r_2)]
=\frac{\mu_0II_0a}{2\pi}
\left(\frac1{r_1}-\frac1{r_2}\right),
\]
整理：
\[
\boxed{F=\frac{2\mu_0II_0}{\pi(4\eta^2-1)}
\simeq0.4\,\mathrm{\mu N}}.
\]
(b) 在恒定电流源约束下
两回路磁相互作用势的改变等于
需要的机械外功。
由电流框两长边位置积分长线磁场的磁通，
原平面构型与翻面构型的互感符号相反，
得到
\[
\boxed{A=\frac{\mu_0aII_0}{\pi}
\ln\frac{2\eta+1}{2\eta-1}
\simeq0.10\,\mathrm{\mu J}}.
\]
\excheck 本题要求机械外功而且\emph{电流恒定}；
不能把固定电流源提供的能量交换忽略。

\electromagproblem{256}{129}
两根很长平行导线半径为 $r$，
轴间距是 $\eta r$，$\eta=20$，
两导线电阻可忽略，
一端由外电阻 $R$ 连接，
另一端接直流电压源。
求两导线之间电磁合力为零时的 $R$。
\exsolution
两线带反号线电荷 $\pm\lambda$，
其横向\emph{静电吸引力}每单位长的模为
$F_e/L=\lambda^2/(2\pi\varepsilon_0d)$，
其中 $d=\eta r$。
两线由远端的电阻与电源闭合成回路，
故两长平行导线中的电流\emph{反向}，
其磁作用为\emph{排斥力}，
与电吸引力相反。磁力模
$F_m/L=\mu_0I^2/(2\pi d)$。
合力为零意味着
$\lambda/I=\sqrt{\mu_0\varepsilon_0}$。
两线每单位长电容
$C'\simeq\pi\varepsilon_0/\ln\eta$，
电压 $V=IR$，
$\lambda=C'V=C'IR$。
于是
\[
\boxed{R=
\frac{\sqrt{\mu_0\varepsilon_0}}{C'}
\simeq\frac{\ln\eta}{\pi}
\sqrt{\frac{\mu_0}{\varepsilon_0}}}.
\]
代 $\sqrt{\mu_0/\varepsilon_0}
\simeq377\,\Omega$，
\[
\boxed{R\simeq0.36\,\mathrm{k\Omega}}.
\]
\excheck 闭合回路使两线电流反向而磁斥，
同时两线所带静态线电荷相反而电吸，
二者才可能通过调节电阻相消。
若电源连接改变两作用方向，结论不再适用。

\electromagproblem{257}{129}
图3.61中的薄半圆截面长导体
通纵向电流 $I$，
另有中心细直线导体
沿\emph{相反方向}通相同电流。
求两导体间的单位长度磁相互作用力。
\exsolution
第225题已得到半圆电流截面
在中心的磁场模
$B_O=\mu_0I/(\pi^2R)$。
中心细线单位长度的磁力
$f=I B_O$，故
\[
\boxed{\frac FL=\frac{\mu_0I^2}{\pi^2R}}.
\]
\excheck 力方向沿半圆形截面对称轴，
由电流方向确定是朝半圆弧还是背离弧；
此处求的是模。

\electromagproblem{258}{129}
图3.71：两根无限长平行导体
通电流 $I_1,I_2$，
第一根为细线，第二根为宽 $b$
的均匀电流带，
最近距离为 $a$。
求两者单位长度磁相互作用力。
\exsolution
把宽带从近边缘 $x=a$
到远边缘 $x=a+b$
分割为并联细电流线，
每个宽度 $\mathrm dx$
中的电流
$\mathrm dI_2=(I_2/b)\mathrm dx$。
细线与其中一段的单位长度磁力
$\mathrm df/L=
\mu_0I_1\,\mathrm dI_2/(2\pi x)$。
积分：
\[
\frac FL=
\frac{\mu_0I_1I_2}{2\pi b}
\int_a^{a+b}\frac{\mathrm dx}{x}
=\boxed{\frac{\mu_0I_1I_2}{2\pi b}
\ln(1+b/a)}.
\]
\excheck 若 $b\ll a$，
对数约为 $b/a$，
恢复两条相距 $a$ 的细线公式。

\electromagproblem{259}{129}
两张无限平行电流片
形成仅在两片之间存在的均匀磁场 $B$，
片外磁场为零。求每片所受单位面积磁力。
\exsolution
对某一片，磁场的法向 Maxwell
应力差与两侧磁能密度差有关。
磁能密度 $w_m=B^2/(2\mu_0)$。
片内侧磁场为 $B$、片外为0，
故片面所受磁压力模
\[
\boxed{p_m=\frac FA=\frac{B^2}{2\mu_0}}.
\]
也可先求每片表面电流密度
$K=B/\mu_0$，
再取\emph{另一片}产生的局部场 $B/2$，
得 $p_m=K(B/2)$。
\excheck 计算自受力时不可直接
用本片两侧磁场的总值作用自身，
否则会多一倍。

\electromagproblem{260}{129--130}
图3.72(a)、(b)、(c)：
无限大电流片放在外加均匀磁场中，
其两侧总磁感应强度分别为
$B_1,B_2$。
求单位面积磁力及图示电流方向。
\exsolution
片面左、右两侧的磁压力分别为
$B_1^2/(2\mu_0)$、
$B_2^2/(2\mu_0)$。
因此沿从左侧指向右侧的正向，
单位面积合力为
\[
\boxed{\frac{F_x}{A}=
\frac{B_1^2-B_2^2}{2\mu_0}}.
\]
若图示 $B_1>B_2$，合力朝右。
电流方向由磁场切向跳变
$\bm K=\bm n\times
(\bm B_{\rm right}-\bm B_{\rm left})/\mu_0$
推定；图3.72的三组箭头对应
\emph{指向纸内}的表面电流。
\excheck 方向须结合图3.72的
\emph{磁场箭头}，
不能只由场强正数 $B_1,B_2$
决定表面电流的有向符号。

\electromagproblem{261}{130}
图3.73电磁泵中，熔融金属
在管道截面积近似为边长
$a=2.0\,\mathrm{cm}$ 的方形区域流动。
跨管道通电流 $I=100\,\mathrm A$，
磁场 $B=0.10\,\mathrm T$
垂直电流和管轴。求泵产生的表压。
\exsolution
电流在磁场中的总安培力
$F=IBa$，沿管道轴向推动金属。
若管截面积为 $a^2$，
泵的表压增量由
$\Delta p\,a^2=F$ 得
\[
\boxed{\Delta p=\frac{IB}{a}}.
\]
代入：
\[
\boxed{\Delta p=
\frac{(100)(0.10)}{0.020}
=500\,\mathrm{Pa}=0.50\,\mathrm{kPa}}.
\]
\excheck 表压正号方向取沿电磁泵
实际推动金属的方向。

\electromagproblem{262}{130}
长薄壁圆柱导体半径为 $R$，
沿轴向通总电流 $I$。
求导体圆柱壁受到的磁压力。
\exsolution
无限长薄壁圆柱内磁场为零，
外部磁场
$B_{\rm out}=\mu_0I/(2\pi R)$。
圆柱壁两侧磁能密度分别为
0 与
$B_{\rm out}^2/(2\mu_0)$。
磁压力（箍缩向内）的模
\[
\boxed{p_m=\frac{B_{\rm out}^2}{2\mu_0}
=\frac{\mu_0I^2}{8\pi^2R^2}}.
\]
\excheck 电流方向不影响压力大小，
因为压力正比于 $I^2$；
磁压力沿径向指向圆柱轴。
% === END 20-magnetostatics-251-262.tex ===

% === BEGIN 20-magnetostatics-263-274.tex blob 68a9810327659ef52437e491494374661e2a3338 ===
% Irodov Part III §3.5 problems 263--274, PDF pp.139--140,
% printed pp.130--131; figures 3.73--3.74 inspected.
\electromagproblem{263}{130}
很长直螺线管单位长度匝数为 $n$，
电流为 $I$。求绕组侧面受到的磁压力。
\exsolution
长螺线管内外磁场分别近似为
$B_{\rm in}=\mu_0nI$ 与 $B_{\rm out}=0$。
侧面磁压力等于两侧磁能密度差：
\[
p_m=\frac{B_{\rm in}^2-B_{\rm out}^2}{2\mu_0}.
\]
故
\[
\boxed{p_m=\frac{\mu_0n^2I^2}{2}}.
\]
\excheck 该压力沿径向\emph{向外}，
螺线管绕组受膨胀张力，
不同于通轴向电流的薄壁圆柱箍缩。

\electromagproblem{264}{130}
长单层螺线管半径 $R$，
单位长度匝数 $n$，
通电流 $I$。
若每一匝导线能够承受的最大拉力为
$F_{\rm lim}$，求断裂临界电流。
\exsolution
单匝在自身之外由其他匝产生的有效磁场
近似是长管内部场的一半：
$B_{\rm eff}\simeq\mu_0nI/2$。
一匝导线单位长度的径向力为
$f=IB_{\rm eff}=\mu_0nI^2/2$。
圆环半径 $R$，用一小段弧的径向受力平衡：
$T\,\mathrm d\theta=fR\,\mathrm d\theta$，
所以
$T=\mu_0nRI^2/2$。
当 $T=F_{\rm lim}$ 时，
\[
\boxed{I_{\rm lim}=
\sqrt{\frac{2F_{\rm lim}}{\mu_0nR}}}.
\]
\excheck $I_{\rm lim}$ 必须带\emph{平方根}，
否则 $2F_{\rm lim}/(\mu_0nR)$
的单位是安培的平方而不是安培。
原书答案纯文本抽取有漏根号风险。

\electromagproblem{265}{130}
平板电极面积 $S$、间距 $d$，
浸没在电阻率为 $\rho$、
以速率 $v$ 平行极板流动的液态导体中。
外磁场 $B$ 在板平面内垂直于流速，
电极以电阻 $R$ 闭合。
求电阻负载上的功率及其最大值。
\exsolution
液体中的动生电场
$\bm E_{\rm mot}=\bm v\times\bm B$，
其跨板距离产生电动势
$\mathcal E_{\rm mot}=vBd$。
液体内部沿板法线的电阻为
$R_{\rm in}=\rho d/S$。
因此负载电流
$I=\mathcal E_{\rm mot}/(R+R_{\rm in})$，
负载功率
\[
\boxed{P_R=
\frac{v^2B^2d^2 R}{(R+\rho d/S)^2}}.
\]
最大功率传输条件
$R=R_{\rm in}$，
于是
\[
\boxed{R_*=\frac{\rho d}{S}},
\quad
\boxed{P_{\max}=
\frac{v^2B^2dS}{4\rho}}.
\]
\excheck 这是液体磁流体发电的
简化模型，忽略涡流与边缘效应；
最大功率工况下内外电阻相等。

\electromagproblem{266}{130}
圆截面铜导线半径
$R=5.0\,\mathrm{mm}$，
沿轴通电流 $I=50\,\mathrm A$。
铜中载流电子数密度
$n=0.9\times10^{23}\,\mathrm{cm^{-3}}$。
求轴线与外表面之间因自磁场产生的横向电势差。
\exsolution
铜线中电子漂移速度模
$v_d=I/(n|e|\pi R^2)$。
导线内部自磁场随半径变化：
$B_\varphi(r)=\mu_0Ir/(2\pi R^2)$。
为平衡漂移电子的径向磁力，
霍尔电场大小
$|E_r|=v_d B_\varphi(r)$。
积分：
\[
|\Delta V|=\int_0^R|E_r(r)|\,\mathrm dr
=\frac{\mu_0Iv_d}{4\pi}.
\]
代漂移速度：
\[
\boxed{|\Delta V|=
\frac{\mu_0I^2}{4\pi^2R^2n|e|}}.
\]
按给定数据，$n=9\times10^{28}\,\mathrm{m^{-3}}$，
数值
\[
\boxed{|\Delta V|\simeq2.2\times10^{-10}\,\mathrm V
=220\,\mathrm{pV}}.
\]
\excheck 本版书末简答数值似为
“$2\,\mathrm{pV}$”，
但其写出的同一解析公式与原题数值
直接代入得到约220pV，
相差约两个数量级；待核查原版物性数值
或单位排印。

\electromagproblem{267}{130--131}
钠导体霍尔测量中，
横向电场
$E_H=5.0\,\mathrm{\mu V/cm}$，
纵向电流密度
$j=200\,\mathrm{A/cm^2}$，
磁感应强度 $B=1.00\,\mathrm T$。
求自由电子数密度及每原子对应自由电子数。
\exsolution
单载流子霍尔平衡
$|e|E_H=|e|v_dB$，
又 $j=n|e|v_d$，
所以
\[
\boxed{n=\frac{jB}{|e|E_H}}.
\]
将 $j=2.0\times10^6\,
\mathrm{A/m^2}$、
$E_H=5.0\times10^{-4}\,\mathrm{V/m}$，
得
\[
\boxed{n\simeq2.5\times10^{28}\,\mathrm{m^{-3}}}.
\]
钠金属原子数密度估计为
$n_{\rm atom}=\rho_{\rm Na}N_A/M_{\rm Na}$，
取 $\rho_{\rm Na}\simeq970\,\mathrm{kg/m^3}$、
$M_{\rm Na}\simeq0.023\,\mathrm{kg/mol}$，
得
$n_{\rm atom}\simeq2.5\times10^{28}\,\mathrm{m^{-3}}$。
因此
$\boxed{n/n_{\rm atom}\simeq1.0}$。
\excheck 原书书末纯文本提取
可能将数密度单位误作 $\mathrm{m^{-1}}$，
正确单位必须是 $\mathrm{m^{-3}}$。

\electromagproblem{268}{131}
铜导体在 $B=100\,\mathrm{mT}$ 下
霍尔测量，横向霍尔电场比纵向电场
小 $\eta=3.1\times10^3$ 倍。
求导电电子的漂移迁移率 $\mu$。
\exsolution
单载流子近似
$E_H=v_dB$，
纵向漂移规律
$v_d=\mu E_\parallel$。
故两场比
$E_H/E_\parallel=\mu B=1/\eta$。
所以
\[
\boxed{\mu=\frac1{\eta B}}.
\]
代 $B=0.100\,\mathrm T$、
$\eta=3.1\times10^3$，
\[
\boxed{\mu\simeq3.2\times10^{-3}\,
\mathrm{m^2/(V\,s)}}.
\]
\excheck \textbf{原书书末数值冲突：}
该版简答似写
$3.2\times10^{-2}\,\mathrm{m^2/(V\,s)}$，
但同一公式 $1/(\eta B)$ 代入0.1T
得到 $3.2\times10^{-3}$，
相差十倍，已保留计算过程。

\electromagproblem{269}{131}
长直导线通电流 $I$，
距它 $r$ 处有磁矩为 $\bm p_m$
的小电流环。
求电流环所受力的模与方向，
当磁矩：(a) 平行直线导体；
(b) 沿径向 $\bm e_r$；
(c) 与当地磁场方向相同。
\exsolution
长线磁场
$\bm B(r)=k_m r^{-1}\bm e_\varphi$，
$k_m=\mu_0I/(2\pi)$。
固定磁矩的力可用
$\bm F=(\bm p_m\cdot\nabla)\bm B$。
(a) $\bm p_m=p_m\bm e_z$，
场与轴向坐标无关，
$\boxed{\bm F_a=0}$。
(b) $\bm p_m=p_m\bm e_r$，
径向导数
$\partial_r\bm B=-k_mr^{-2}\bm e_\varphi$，
\[
\boxed{\bm F_b=
-\frac{\mu_0Ip_m}{2\pi r^2}\bm e_\varphi}.
\]
(c) $\bm p_m=p_m\bm e_\varphi$，
沿方位角方向移动时
$\partial_\varphi\bm e_\varphi=-\bm e_r$，
定向导数
$(p_m/r)\partial_\varphi\bm B
=-p_mk_m r^{-2}\bm e_r$，
故
\[
\boxed{\bm F_c=
-\frac{\mu_0Ip_m}{2\pi r^2}\bm e_r}.
\]
\excheck (b) 的力与当地 $\bm B$ 反向；
(c) 指向导线轴；
符号随磁矩或电流反向而改变。

\electromagproblem{270}{131}
半径 $R$、电流 $I$ 的圆环轴线上
距中心 $x$ 处放一只小线圈，
磁矩 $p_m$ 平行圆环轴。
求线圈所受磁力模。
\exsolution
大圆环轴向场
$B(x)=\mu_0IR^2/
[2(R^2+x^2)^{3/2}]$。
小磁偶极子在非均匀场中受力
$F_x=p_m\,\mathrm dB/\mathrm dx$，
其模
\[
|F_x|=p_m\frac{3\mu_0IR^2|x|}
{2(R^2+x^2)^{5/2}}.
\]
故
\[
\boxed{F=
\frac{3\mu_0IR^2p_m|x|}
{2(R^2+x^2)^{5/2}}}.
\]
\excheck $x=0$ 因对称性受力为零，
远场 $x\gg R$ 时按 $x^{-4}$ 衰减。

\electromagproblem{271}{131}
两个很小的同轴线圈磁矩为
$p_{m1}=4.0\,\mathrm{mA\,m^2}$、
$p_{m2}=6.0\,\mathrm{mA\,m^2}$，
间距 $\ell=20\,\mathrm{cm}$
远大于线圈大小。
求相互作用力大小。
\exsolution
磁偶极子1的轴上远场
$B(\ell)=\mu_0p_{m1}/(2\pi\ell^3)$。
第二个同轴磁偶极子的力模
\[
F=p_{m2}\left|\frac{\mathrm dB}{\mathrm d\ell}\right|
=\frac{3\mu_0p_{m1}p_{m2}}{2\pi\ell^4}.
\]
代数值：
\[
\boxed{F\simeq9.0\times10^{-9}\,\mathrm N=9\,\mathrm{nN}}.
\]
\excheck 两磁矩同向时吸引，
反向时排斥；题目未给相对方向时
所列为力的\emph{大小}。

\electromagproblem{272}{131}
薄圆盘形永磁体沿轴磁化，
盘半径 $R=1.0\,\mathrm{cm}$。
在轴线上距离 $x=10\,\mathrm{cm}$
处测得磁场 $B=30\,\mathrm{\mu T}$。
估算沿圆盘边缘流动的等效分子电流 $I'$。
\exsolution
薄轴向磁化圆盘的边缘束缚电流
可等效为半径 $R$ 的电流环，
轴上磁场
$B(x)=\mu_0I'R^2/
[2(R^2+x^2)^{3/2}]$。
当 $x\gg R$，
$B\simeq\mu_0I'R^2/(2x^3)$，
故
\[
\boxed{I'\simeq\frac{2Bx^3}{\mu_0R^2}}.
\]
代入 SI 数值：
\[
\boxed{I'\simeq4.8\times10^2\,\mathrm A
\approx0.5\,\mathrm{kA}}.
\]
\excheck 这是\emph{磁化分子束缚电流}的
宏观等效，不是外电源实际输送的自由电流。

\electromagproblem{273}{131}
均匀各向同性磁介质和真空
之间有平界面，介质相对磁导率 $\mu$。
真空侧磁感应强度模为 $B$，
与界面法线夹角 $\alpha$。
求介质内近界面场强模 $B'$。
\exsolution
若界面无自由面电流，
磁感应强度的法向分量连续：
$B'_n=B_n=B\cos\alpha$。
磁场强度的切向分量连续：
$H'_t=H_t=B\sin\alpha/\mu_0$。
因此介质内
$B'_t=\mu\mu_0H'_t=\mu B\sin\alpha$。
模的平方为
$B'^2=B_n'^2+B_t'^2$，
所以
\[
\boxed{B'=B\sqrt{\cos^2\alpha+
\mu^2\sin^2\alpha}}.
\]
\excheck $\alpha=0$ 时 $B'=B$，
$\alpha=\pi/2$ 时 $B'=\mu B$，
分别是法向与切向边界条件的直接极限。

\electromagproblem{274}{131--132}
图3.74：磁介质相对磁导率 $\mu$，
与真空以平面分界。
真空侧磁感应强度大小 $B$，
与法线 $\bm n$ 夹角 $\theta$。
(a) 求以界面上一点为球心、半径 $R$
的球面上 $\bm H$ 的通量；
(b) 求图示横跨界面的边长 $\ell$
的正方形闭路上 $\bm B$ 的环流。
\exsolution
法向 $B_n=B\cos\theta$ 连续，
所以
$H_{{\rm vac},n}=B\cos\theta/\mu_0$，
$H_{{\rm mag},n}=B\cos\theta/(\mu\mu_0)$。
(a) 两半球的法向投影面积均为
$\pi R^2$，法线对整个闭合球面
方向相反，净通量
\[
\boxed{\oint_S\bm H\cdot\mathrm d\bm S
=\frac{\pi R^2B\cos\theta}{\mu_0}
\left(1-\frac1\mu\right)}.
\]
(b) 磁场切向分量在两侧为
$B_{{\rm vac},t}=B\sin\theta$、
$B_{{\rm mag},t}=\mu B\sin\theta$。
沿图示正向闭合方框，平行边贡献相减，
垂直边因边长趋零或场分量相消，
得到
\[
\boxed{\oint_\Gamma\bm B\cdot\mathrm d\bm\ell
=(1-\mu)B\ell\sin\theta}.
\]
\excheck 不能把 $\bm B$ 的环流
与 $\bm H$ 的环流混淆：
前者可因束缚电流而非零。
% === END 20-magnetostatics-263-274.tex ===

% === BEGIN 20-magnetostatics-275-287.tex blob aafa1fecda8ef208f54f636e355802baf3e22254 ===
% Irodov Part III §3.5 last 275--287; PDF pp.141--143,
% printed pp.132--134. Graph of original Fig 3.76 consulted.
\electromagproblem{275}{132}
一根很长的圆柱形顺磁导线
流轴向总电流 $I$，
均匀磁化率为 $\chi$。
求：(a) 导线表面束缚分子电流总量；
(b) 体内束缚分子电流总量；
并说明相对流向。
\exsolution
导线半径记为 $R$。
自由电流密度
$j=I/(\pi R^2)$，
线内
$H_\varphi(r)=jr/2=Ir/(2\pi R^2)$。
磁化强度
$J_\varphi(r)=\chi H_\varphi(r)
=\chi I r/(2\pi R^2)$。
体束缚电流密度
$\bm j_b=\nabla\times\bm J$，
其轴向分量
\[
j_{b,z}=\frac1r\frac{\mathrm d}{\mathrm dr}
(rJ_\varphi)
=\frac{\chi I}{\pi R^2}.
\]
截面积分体束缚电流
\[
\boxed{I'_{\rm vol}=\chi I}.
\]
表面束缚电流密度
$\bm K_b=\bm J\times\bm n$，
在 $r=R$，
$\bm J\parallel\bm e_\varphi$、
$\bm n=\bm e_r$，
故 $\bm K_b\parallel-\bm e_z$，
沿周长积分模为
\[
\boxed{I'_{\rm surf}=
2\pi R\,J_\varphi(R)=\chi I}.
\]
\excheck 两者\emph{大小相等、方向相反}，
导线及外侧整体的束缚电流净量为零。

\electromagproblem{276}{132}
图3.75：很长的通电螺线管
从中部 $O$ 开始，轴线 $x>0$
的半段填充相对磁导率为 $\mu$ 的介质，
$x<0$ 为真空。
画出轴线上 $B(x)$、$H(x)$、
磁化强度 $J(x)$ 的大致分布。
\exsolution
远离分界面，螺线管的自由表面电流
等效密度为 $K=nI$。
真空一侧（$x\ll-R$）
\[
\boxed{H_-\simeq nI,\qquad
B_-\simeq\mu_0nI,\qquad J_-=0}.
\]
介质一侧（$x\gg R$）
\[
\boxed{H_+\simeq nI,\qquad
B_+\simeq\mu\mu_0nI,\qquad
J_+\simeq(\mu-1)nI}.
\]
在 $O$ 附近，由端部和磁荷退极化效应，
$B$、$H$ 的真实变化不能只以
两边常值拼接；沿轴无自由磁荷，
磁感应强度的法向分量连续，
故 $B(x)$ 应平滑连接两侧平台。
磁化强度在介质外为零，
进入介质后建立并趋于平台。
\excheck 此处说明的是原书要求的
\emph{定性图形}，不是磁介质有限端部精确边值解。
若只把 $\mu$ 作为乘数应用到整个轴上，
会错误地改变真空侧远处的场。

\electromagproblem{277}{132}
无限长直导线通电流 $I$，
恰位于两种非导电磁介质的平面交界线上，
相对磁导率分别为 $\mu_1,\mu_2$。
磁力线为以导线为中心的圆周，
求距导线 $r$ 处磁感应强度模。
\exsolution
在同一环路上，$\bm B$
沿界面切向的法向跨界条件以及
题设圆形磁力线使其模在两半圆上相同。
相应磁场强度分别为
$H_1=B/(\mu_0\mu_1)$、
$H_2=B/(\mu_0\mu_2)$。
半径为 $r$ 的整圆安培环路由两半圆组成，
\[
I=\oint\bm H\cdot\mathrm d\bm\ell
=\pi r\frac B{\mu_0\mu_1}
+\pi r\frac B{\mu_0\mu_2}.
\]
解得
\[
\boxed{B(r)=\frac{\mu_0I\,\mu_1\mu_2}
{\pi r(\mu_1+\mu_2)}}.
\]
\excheck $\mu_1=\mu_2=1$ 时
恢复真空直导线场
$B=\mu_0I/(2\pi r)$。

\electromagproblem{278}{132}
半径为 $R$ 的环形电流位于
真空与磁介质的\emph{平面边界内}，
磁介质相对磁导率 $\mu$。
若无磁介质时场为 $\bm B_0$，
求该构型在电流环轴上及其余空间的磁感应强度。
\exsolution
把电流环所在平面作为介质分界面，
两侧均是无自由体电流区域。
对各个毕奥—萨伐尔自由电流元产生的
磁矢势使用线性材料边界条件，
可将本几何中的总场写成
原均匀真空问题的固定比例：
\[
\bm B(\bm r)=\frac{2\mu}{1+\mu}\bm B_0(\bm r).
\]
因此轴上及其它适用的场点：
\[
\boxed{B=\frac{2\mu}{1+\mu}B_0}.
\]
\excheck $\mu=1$ 还原 $\bm B_0$；
$\mu\to\infty$ 比例趋于2。
该结果依赖环电流\emph{严格位于两介质平界面}
的特殊几何，而非任意介质分布都适用。

\electromagproblem{279}{132}
相对磁导率为 $\mu$ 的
均匀各向同性磁介质球
置于外加均匀磁场 $\bm B_0$ 中，
球内均匀磁化。
已知均匀磁化球的退磁场
$\bm H_d=-\bm J/3$。
求球内磁感应强度。
\exsolution
外加磁场强度
$\bm H_0=\bm B_0/\mu_0$，
球内
$\bm H=\bm H_0-\bm J/3$。
线性材料磁化
$\bm J=(\mu-1)\bm H$，
因此
\[
\bm H=\bm H_0-\frac{\mu-1}{3}\bm H,
\quad
\bm H=\frac3{\mu+2}\bm H_0.
\]
材料内
$\bm B=\mu\mu_0\bm H$，所以
\[
\boxed{\bm B=
\frac{3\mu}{\mu+2}\bm B_0}.
\]
\excheck $\mu\to1$ 时系数为1；
$\mu\to\infty$ 时系数趋于3。

\electromagproblem{280}{132}
长 $\ell=15\,\mathrm{cm}$ 的
圆柱形永久磁铁上均匀绕
$N=300$ 匝线圈。
通电流 $I=3.0\,\mathrm A$ 后
磁铁外磁场恰好消失，
求磁性材料的矫顽力 $H_c$。
\exsolution
线圈在长圆柱轴上产生磁场强度
$H_{\rm coil}\simeq nI=NI/\ell$。
要抵消永久磁化所产生的外场，
需施加相反方向的矫顽磁场。
故矫顽力模为
\[
\boxed{H_c=\frac{NI}{\ell}}.
\]
代数值：
\[
\boxed{H_c=\frac{(300)(3.0)}{0.15}
=6.0\times10^3\,\mathrm{A/m}
=6.0\,\mathrm{kA/m}}.
\]
\excheck 题目测量使用的是磁场强度
$\bm H$ 而不是磁感应强度 $\bm B$，
单位为 $\mathrm{A/m}$。

\electromagproblem{281}{132}
永久磁铁制成平均直径
$d=20\,\mathrm{cm}$ 的闭合圆环，
环上留宽
$b=2.0\,\mathrm{mm}$ 的小气隙。
气隙内磁感应强度
$B=40\,\mathrm{mT}$。
忽略边缘漏磁，求磁铁内部磁场强度模。
\exsolution
没有自由绕组电流，
沿闭合环路有
$\oint\bm H\cdot\mathrm d\bm\ell=0$。
设铁磁体中的磁场强度为 $H_m$，
气隙中
$H_g=B/\mu_0$。
忽略气隙长度对主圆周的微小影响：
\[
H_m(\pi d-b)+\frac B{\mu_0}b=0.
\]
因此
\[
\boxed{|H_m|=
\frac{bB}{\mu_0(\pi d-b)}
\simeq\frac{bB}{\mu_0\pi d}}.
\]
代入：
\[
\boxed{|H_m|\simeq1.0\times10^2\,
\mathrm{A/m}=0.10\,\mathrm{kA/m}}.
\]
\excheck 磁铁内部的 $\bm H_m$
与气隙场方向相反；
此处报磁场强度的模。

\electromagproblem{282}{132--133}
平均半径
$R=250\,\mathrm{mm}$ 的环形铁芯，
绕有 $N=1000$ 匝，
电流 $I=0.85\,\mathrm A$。
铁芯切有气隙
$b=1.00\,\mathrm{mm}$，
气隙磁感应强度
$B=0.75\,\mathrm T$。
忽略漏磁，求该工况下铁的相对磁导率 $\mu$。
\exsolution
沿环路应用安培定律：
\[
NI=H_{\rm iron}(2\pi R-b)+H_{\rm gap}b.
\]
两处磁通守恒且截面相同，
$B_{\rm iron}\simeq B_{\rm gap}=B$。
因此
$H_{\rm iron}=B/(\mu\mu_0)$、
$H_{\rm gap}=B/\mu_0$。
解出
\[
\boxed{\mu=
\frac{B(2\pi R-b)}
{\mu_0NI-Bb}}.
\]
在 $b\ll R$ 下近似
$\mu\simeq2\pi RB/(\mu_0NI-Bb)$。
代入：
\[
\boxed{\mu\simeq3.7\times10^3}.
\]
\excheck 分母 $\mu_0NI-Bb$ 必须为正，
否则所假设的同向磁化与场分布
不适用于题设工况。

\electromagproblem{283}{133}
原书图3.76给出了商业纯铁的初始磁化曲线
$B(H)$。由图推绘相对磁导率 $\mu(H)$，
并估计最大值与对应磁场强度。
\exsolution
在各个图上的非零 $H$ 点，
相对磁导率按定义为
\[
\boxed{\mu(H)=\frac{B(H)}{\mu_0H}}.
\]
因此应逐点用纵坐标 $B$ 除以
横坐标 $H$（以 SI
$\mathrm{A/m}$ 代入），再把所得点作图。
在 $H$ 很小时，$B$ 起初随 $H$ 迅速增加，
$\mu$ 上升；到饱和区
$B$ 增长缓慢，而分母 $H$ 增大，
故 $\mu$ 下降。
从原书图3.76读数，
最大值大约位于
\[
\boxed{H_{\max}\simeq0.06\,\mathrm{kA/m}
=60\,\mathrm{A/m}},
\]
该处磁感应强度约 $0.75\,\mathrm T$，
所以
\[
\boxed{\mu_{\max}\approx
\frac{0.75}{(4\pi\times10^{-7})(60)}
\simeq1.0\times10^4}.
\]
\excheck 此题采用\emph{从印刷图读数}
的近似值，不具有精密物性测量的有效数字；
$\mu(H)$ 是非线性材料的割线磁导率。

\electromagproblem{284}{133}
图3.76同一铁磁材料，
平均直径
$d=50\,\mathrm{cm}$ 的细环形铁芯
绕 $N=800$ 匝，
电流 $I=3.0\,\mathrm A$，
铁芯气隙宽
$b=2.0\,\mathrm{mm}$。
忽略漏磁，结合原图求工作点 $\mu$。
\exsolution
沿磁路
$NI=H(\pi d-b)+Hb_g$，
其中气隙内
$H_g=B/\mu_0$，
因此
\[
B=\frac{\mu_0NI}{b}
-\frac{\mu_0(\pi d-b)}{b}\,H.
\]
代入数据并把 $H$ 的横坐标写作
$\mathrm{kA/m}$：
\[
\boxed{B(\mathrm T)\simeq
1.51-0.987H(\mathrm{kA/m})}.
\]
与原书图3.76的非线性磁化曲线求交，
读得
$H\simeq0.26\,\mathrm{kA/m}$、
$B\simeq1.25\,\mathrm T$。
相对磁导率
\[
\boxed{\mu=\frac B{\mu_0H}
\simeq\frac{1.25}
{(4\pi\times10^{-7})(260)}
\simeq3.8\times10^3\approx4\times10^3}.
\]
\excheck \textbf{书后答案数量级冲突：}
书后简答似给“$4\times10^2$”，
但使用同一书后读数 $B\approx1.25$T、
$H\approx0.26$kA/m，
必得 $\sim4\times10^3$，
相差10倍；本稿按 SI 数值计算。

\electromagproblem{285}{133}
横截面积 $S$、磁化率为 $\chi$
的细长顺磁棒，
沿通电线圈轴向放置。
一端位于线圈中心、局部磁感应强度为 $B$，
另一端位于磁场可以忽略的远处。
求线圈对棒的磁力。
\exsolution
弱顺磁体的磁化
$J=\chi H\simeq\chi B/\mu_0$。
恒定线圈电流下，
棒单位体积的可用磁势能近似
$u=-\chi B^2/(2\mu_0)$。
棒沿轴微量进入线圈，
相当于有体积 $S\,\mathrm dx$
从无场处移到场强 $B$ 处。
有效能下降
$\mathrm dU\simeq
-[\chi B^2/(2\mu_0)]S\,\mathrm dx$，
吸入力模为
\[
\boxed{F\simeq\frac{\chi SB^2}{2\mu_0}}.
\]
\excheck 方向为将顺磁棒\emph{吸入}
强磁场一侧；
该式只适用于弱磁化与细棒近似。

\electromagproblem{286}{133--134}
图3.77中体积
$V=41\,\mathrm{mm^3}$
的顺磁小球位于电磁铁极靴轴线上。
磁场随高度变化为
$B(x)=B_0e^{-a x^2}$，
$B_0=1.50\,\mathrm T$、
$a=100\,\mathrm{m^{-2}}$。
求 (a) 磁吸引力最大的高度；
(b) 若 $F_{\max}=160\,\mathrm{\mu N}$，
求材料磁化率 $\chi$。
\exsolution
弱顺磁小球的磁势能近似为
$U(x)=-\chi VB(x)^2/(2\mu_0)$。
磁力沿 $x$ 方向
$F_x=-\mathrm dU/\mathrm dx
=[\chi V/(2\mu_0)]\,\mathrm d(B^2)/\mathrm dx$。
因为
$B^2=B_0^2e^{-2ax^2}$，
吸引力模
\[
|F(x)|=\frac{2\chi Va B_0^2}{\mu_0}
x\,e^{-2ax^2},\quad x>0.
\]
对 $x e^{-2ax^2}$ 求导，
极值条件
$1-4ax^2=0$，于是
\[
\boxed{x_{\max}=\frac1{2\sqrt a}
=0.050\,\mathrm m=5.0\,\mathrm{cm}}.
\]
在该处
$|F_{\max}|=
\chi VB_0^2\sqrt a\,e^{-1/2}/\mu_0$，
所以
\[
\boxed{\chi=
\frac{\mu_0F_{\max}e^{1/2}}
{VB_0^2\sqrt a}
\simeq3.6\times10^{-4}}.
\]
\excheck 球体积换算
$41\,\mathrm{mm^3}=41\times10^{-9}\,\mathrm{m^3}$。
书后纯文本提取的 $x_m$ 可能丢平方根；
从极值方程所得长度尺度必须是 $1/\sqrt a$。

\electromagproblem{287}{134}
一颗体积为 $V$、
顺磁磁化率为 $\chi$ 的小球，
从线圈轴上磁感应强度为 $B$ 的位置，
被缓慢移动到场可以忽略的远处。
求外力所做功。
\exsolution
弱线性顺磁球在给定外场中的
等效磁势能（保持线圈电流不变）
为
$U=-\chi VB^2/(2\mu_0)$。
初态
$U_i=-\chi VB^2/(2\mu_0)$，
末态
$U_f\simeq0$。
准静态移走小球的外力功
$A_{\rm ext}=U_f-U_i$：
\[
\boxed{A_{\rm ext}\simeq
\frac{\chi VB^2}{2\mu_0}}.
\]
\excheck 对顺磁材料 $\chi>0$，
外界需克服吸引力做正功；
若改为抗磁材料 $\chi<0$，
力方向与功号将相反。
% === END 20-magnetostatics-275-287.tex ===
